\section{从 Actor-Critic 到纯 Value-Based 方法}

前几讲介绍了 actor-critic 方法：同时维护策略（actor）和价值函数（critic）。Q-Learning 走得更远——\textbf{完全去掉显式策略}，只学习 Q 函数，然后从 Q 函数隐式地派生策略。

\begin{knowledgebox}{Value-Based vs. Actor-Critic}
\begin{itemize}
    \item \textbf{Actor-Critic}：同时学习策略 $\pi_\theta$ 和价值函数 $\hat{V}$ 或 $\hat{Q}$。
    \item \textbf{Value-Based (Q-Learning)}：只学习最优 Q 函数 $Q^*$，策略隐式为 $\pi(s) = \arg\max_a Q^*(s, a)$。
\end{itemize}
\end{knowledgebox}

\subsection{本章小结}

Q-Learning 是一种纯 value-based 方法，不需要显式的策略网络。

